\documentclass[11pt]{article}
\usepackage{amsmath,amssymb}
\usepackage{geometry}
\geometry{margin=1in}

\title{Bias-Aware Isotropic Transformer:\\
Mean Forwarding and Spike Whitening at the Embedding Interface}
\author{}
\date{}

\begin{document}
\maketitle

\section{Motivation}

Empirical analysis of large language model training reveals a consistent geometric pattern:

\begin{itemize}
    \item The \textbf{embedding activation} exhibits a large, statistically significant mean vector.
    \item It also exhibits strong spectral anisotropy: a small number of singular directions dominate the variance.
    \item In contrast, gradient outputs do not exhibit a comparable mean bias.
\end{itemize}

This suggests that much of the anisotropy observed in later transformer layers may originate from the embedding layer itself.

Importantly, the bias and spike are not errors in the data. They reflect:

\begin{itemize}
    \item Frequency structure of natural language.
    \item Shared grammatical components.
    \item Common background statistics of token distributions.
\end{itemize}

However, when propagated unmodified through deep residual stacks and adaptive optimizers, this structured bias can:

\begin{itemize}
    \item Inflate second-moment statistics.
    \item Distort attention geometry.
    \item Dominate dynamic range in low-bit quantization.
\end{itemize}

The key design question is therefore:

\begin{quote}
How can we preserve structural bias while preventing it from corrupting semantic processing?
\end{quote}

\section{Design Principle}

Two empirical findings guide the design:

\begin{enumerate}
    \item The dominant outliers in activations are mean-dominated.
    \item The spike subspace encodes structured variance, not pure noise.
\end{enumerate}

Therefore:

\begin{itemize}
    \item We should \textbf{separate the mean bias}, but not delete it.
    \item We should \textbf{whiten the spike subspace}, not remove it.
\end{itemize}

This leads to a dual-path design:

\begin{itemize}
    \item A \textbf{semantic path} that is mean-free and spectrally flattened.
    \item A \textbf{structure path} that forwards the mean to the output layer.
\end{itemize}

\section{Mathematical Formulation}

Let the embedding activation for a batch be:

\[
X \in \mathbb{R}^{N \times H}.
\]

\subsection{Step 1: Mean Decomposition}

Compute the batch mean:

\[
\mu = \frac{1}{N} \mathbf{1}^\top X \in \mathbb{R}^{H}.
\]

Decompose:

\[
X = \mathbf{1}\mu^\top + X_c,
\]
where $X_c$ is mean-centered.

\subsection{Step 2: Spike Whitening}

Compute SVD of centered activation:

\[
X_c = U \Sigma V^\top.
\]

Let $V_k$ be the top-$k$ spike directions (e.g., top $1\%$).

Instead of removing them, whiten them:

\[
\Sigma_k \rightarrow I.
\]

Define whitening operator:

\[
W = V_k \Sigma_k^{-1} V_k^\top + (I - V_k V_k^\top).
\]

Whitened activation:

\[
X_{\text{iso}} = X_c W.
\]

This preserves directions while equalizing their variance.

\subsection{Step 3: Transformer Processing}

Feed only the isotropic semantic stream into transformer layers:

\[
H_0 = X_{\text{iso}}.
\]

\[
H_{\ell+1} = H_\ell + \mathcal{F}_\ell(H_\ell).
\]

Because $H_0$ is zero-mean and spectrally flattened:

\begin{itemize}
    \item Attention is not dominated by frequency directions.
    \item Optimizer second moments are better conditioned.
\end{itemize}

\subsection{Step 4: Mean Forwarding}

The accumulated structural mean is forwarded separately:

\[
x_{\text{final}} = H_L + \mu.
\]

\[
\text{Logits} = \text{Unembed}(x_{\text{final}}).
\]

Thus:

\begin{itemize}
    \item Semantic reasoning occurs in an isotropic space.
    \item Frequency prior is preserved for prediction.
\end{itemize}

\section{Interpretation}

This architecture separates:

\begin{itemize}
    \item \textbf{Carrier Wave (Structure):} mean bias $\mu$.
    \item \textbf{Signal (Semantics):} zero-mean, whitened activation.
\end{itemize}

It avoids:

\begin{itemize}
    \item Deleting spike information (as in hard SVD truncation).
    \item Allowing frequency-dominated directions to distort geometry.
\end{itemize}

\section{Expected Benefits}

\begin{itemize}
    \item Reduced dynamic-range inflation for low-bit quantization.
    \item Improved optimizer conditioning.
    \item Preservation of syntactic and frequency structure.
    \item Cleaner semantic geometry in deep layers.
\end{itemize}

\section{Conclusion}

The empirical evidence suggests that anisotropy and mean bias originate at the embedding interface.
Rather than treating anisotropy as a property of deep layers, we propose handling it at its source:

\begin{itemize}
    \item Separate mean bias.
    \item Whiten spike variance.
    \item Preserve structure via forwarding.
\end{itemize}

This provides a principled architecture that aligns model geometry with the statistical structure of language, without discarding information.

\end{document}